\section{Introduction}\label{sec:intro}

\subsection{A chaotic, high-dimensional, memory-bound object}
Numerical weather prediction (NWP) integrates a discretised set of partial differential equations
for momentum, mass, thermodynamics and water species forward in time. The state is large: a
0.25$^\circ$ global model carries of order $10^{8}$ grid points, a 3 km regional ensemble carrying
$10^{7}$ points per member. The dynamics are chaotic. Two trajectories that agree to the last bit
diverge at a rate set by the leading Lyapunov exponent, and the practical forecast horizon is the
time it takes an initial error to grow to the scale of the analysis uncertainty. For the
Lorenz--96 system used throughout this report as a tractable surrogate, the leading exponent is
$\lambda_1 \approx 0.34$ day$^{-1}$ and the error-doubling time is \ACDoubling\ days
(Fig.~\ref{fig:req}(b,c)) --- numbers that are, by construction, in the same range as operational
atmospheric error growth \cite{lorenz1963,lorenz1996,kalnay2003}.

Uncertainty is handled by running an ensemble. A deterministic forecast answers ``what will the
weather be''; an ensemble of 30--100 members, each with perturbed initial conditions and
stochastic physics, answers ``how likely is it''. Ensembles are what make NWP useful for
decision-making, and they are also what make it expensive: cost scales linearly in members, and
the linear scaling is not amortised by any cache or reuse because each member is an independent
trajectory that nevertheless shares the same grid, coefficients and boundary data.

Two further facts shape the machine this report specifies.

\textbf{Arithmetic intensity is low and fixed.} A 3D stencil that reads $b$ bytes per cell and
performs $f$ flops per cell has an intensity of $f/b$ that cannot be improved by tiling beyond the
size of the on-chip memory. For the 24-variable physics package used here, $f \approx 4{,}540$
flops per cell per member and $b = 192$ bytes, i.e. 23.6 flop/byte at the HBM boundary. That is
one of the lowest intensities in high-performance computing, and it is why the kernel is a memory
problem with arithmetic attached rather than the other way around.

\textbf{The result must be reproducible and auditable.} Operational centres must be able to
reproduce a forecast bit-for-bit for post-event analysis, and dispersion or aviation-hazard
products are increasingly treated as safety-relevant. A machine whose numerical path contains
atomics or dynamic scheduling cannot make that guarantee cheaply.

\subsection{Why a GPU is a poor fit at the kernel level}
Graphics processors became the default accelerator for NWP because they offer the highest
floating-point throughput per dollar. On the kernels that dominate an atmospheric model, however,
three structural costs remain.

\textbf{Issue and data-movement overhead.} A CUDA core retires one fp32 FMA per lane per cycle at
best, and a warp instruction carries address arithmetic, predicate handling and register-file
traffic around it. Published energy numbers put a delivered fp32 FLOP on a data-centre GPU at
several picojoules once the memory hierarchy is included; the arithmetic itself is a small
fraction of that \cite{horowitz2014}. On a stencil with 23 flop/byte, the memory hierarchy is
exercised on every FLOP.

\textbf{Kernel-boundary re-fetch.} Because a GPU kernel is launched per operator, the state is
written to and read back from HBM at each kernel boundary. Fusing an entire physics package into
one kernel is possible in principle and hard in practice; the reference model assumes a re-fetch
factor of 2.5, i.e. the 25.8 GB state is streamed 2.5 times per time step rather than once.

\textbf{Precision inflexibility.} A GPU is fastest in fp16 or fp8 on tensor cores, formats that a
chaotic model cannot use for its state, and slowest in the fp64 that some centres still use for
conservation-sensitive tracers. The useful middle ground --- fp32 dynamics with fp16 or bf16
physics, chosen per grid cell from the local ensemble spread --- has no hardware expression on a
general-purpose SIMT machine.

None of this makes GPUs bad. It makes them general. The question this report asks is what a
machine looks like when the ensemble dimension, the halo structure and the error-growth physics are
hardware primitives rather than software conventions.

\subsection{Thesis}
\ac{1} is built on three principles, each of which is falsifiable and each of which is tested in
Sections~\ref{sec:results} and~\ref{sec:validation}.

\begin{enumerate}[leftmargin=1.4em,itemsep=1pt]
\item \textbf{Lanes carry members, not matrix rows.} 256 lanes per tile, 64 tiles, 16,384 lanes in
      total. A stencil coefficient is broadcast to all lanes; a grid neighbour is exchanged through
      a dedicated six-port fabric; an ensemble member never leaves its lane. Utilisation on the
      reference kernels is 71.1\% rather than the 20--30\% typical of a stencil on a GPU.
\item \textbf{Bytes are the budget.} The machine is designed so that a 24-variable, 3-time-level
      tile lives in 256 KB of SRAM, a fused pass streams the 25.8 GB state exactly once per step,
      and the halo traffic that a GPU sends through L2/L3 never leaves the register file. 56\% of
      the energy is HBM traffic; the design is honest about that and optimises it first.
\item \textbf{Precision follows predictability.} Where the ensemble spread is large, the forecast
      error is dominated by chaotic uncertainty, not by the 11-bit mantissa of fp16. A per-cell
      governor exploits this: 58\% of physics work by design (84\% by measurement) runs in fp16
      with no detectable skill loss, while dynamics stays fp32.
\end{enumerate}

\subsection{Contributions}
\begin{itemize}[leftmargin=1.4em,itemsep=1pt]
\item An architecture and ISA for ensemble-native atmospheric computation, with a cycle-level
      model of its tile pipeline and a 64-tile chip model (Sections~\ref{sec:arch}--\ref{sec:micro}).
\item A precision-governor design grounded in a real cycled ensemble experiment, including the
      negative result that fp8 is not usable for the state (Section~\ref{sec:precision}).
\item A calibrated pre-silicon evaluation with ablations, a 300-sample Monte-Carlo sensitivity
      study, and an explicit ``idealised GPU'' bound (Section~\ref{sec:results}).
\item A 3 nm GAA process integration with backside power delivery, hybrid-bonded L3, 2.5D HBM
      packaging, yield/cost model and qualification plan (Sections~\ref{sec:process}--\ref{sec:mfg}).
\item A complete software stack: compiler flow, driver ABI, Python/Node SDKs and a cross-platform
      desktop application that calls the device through one-line functions (Section~\ref{sec:software}).
\end{itemize}

\subsection{Scope and honesty}
This is a design study. There is no silicon, no RTL netlist and no place-and-route database behind
the numbers. The evaluation is a transparent performance model whose constants are either taken
from the literature or stated as assumptions, and whose output is reported with an interval rather
than a point. Where the model is weak --- the DA comparison, the GPU efficiency assumption, the
absence of MPI scale-out --- the report says so. The value of the exercise is the architecture and
the reasoning, not the fourth digit of a projected number.

\subsection{Organisation}
Section~\ref{sec:workload} characterises the workload and the benchmark.
Section~\ref{sec:arch} describes the chip, Section~\ref{sec:micro} the tile and its ISA.
Section~\ref{sec:precision} develops the precision and energy theory.
Section~\ref{sec:results} reports performance and energy, Section~\ref{sec:validation} the
numerical validation. Sections~\ref{sec:process} and~\ref{sec:mfg} cover silicon and
manufacturing. Section~\ref{sec:software} the software stack, Section~\ref{sec:deployment}
deployment and limitations, Section~\ref{sec:roadmap} the roadmap.
